\documentclass[11pt]{article}
\usepackage{amsmath,amssymb,mathtools,bm}
\usepackage[margin=1in]{geometry}
\usepackage{booktabs}
\usepackage{hyperref}

\title{TorchTEM Layer Equations}
\author{tspi}
\date{}

\begin{document}
\maketitle

\section{Purpose}
This document records the mathematical form of implemented layers in
\texttt{torchtem}.  

\section{Notation}
We use the following notation:
\begin{align}
  \psi(\mathbf{r}) &\in \mathbb{C} && \text{complex wave in real space,} \\
  \Psi(\mathbf{k}) &\in \mathbb{C} && \text{complex wave in reciprocal space,} \\
  \mathbf{r} &= (x,y), \\
  \mathbf{k} &= (k_x,k_y), \\
  \lambda &= \lambda(E) && \text{relativistic electron wavelength,} \\
  \sigma &= \sigma(E) && \text{electron interaction parameter,} \\
  V_j(\mathbf{r}) &&& \text{projected electrostatic potential of slice } j.
\end{align}
The forward and inverse Fourier transforms in the current implementation use
orthonormal FFT normalization:
\begin{align}
  \Psi(\mathbf{k}) &= \mathcal{F}\{\psi(\mathbf{r})\}, \\
  \psi(\mathbf{r}) &= \mathcal{F}^{-1}\{\Psi(\mathbf{k})\}.
\end{align}

\section{Wave Containers and Builders}
\subsection{Waves}
The Torch-native \texttt{Waves} container stores an array
\begin{align}
  \psi \in \mathbb{C}^{n_1 \times \cdots \times n_m \times N_y \times N_x},
\end{align}
where the final two indices are spatial dimensions and the leading indices are optional
ensemble dimensions.

The real-space intensity and phase observables are
\begin{align}
  I(\mathbf{r}) &= |\psi(\mathbf{r})|^2, \\
  \phi(\mathbf{r}) &= \arg \psi(\mathbf{r}).
\end{align}

Reciprocal-space conversion is defined by
\begin{align}
  \Psi(\mathbf{k}) &= \mathcal{F}\{\psi(\mathbf{r})\},
\end{align}
and real-space reconstruction from reciprocal representation is
\begin{align}
  \psi(\mathbf{r}) &= \mathcal{F}^{-1}\{\Psi(\mathbf{k})\}.
\end{align}

Two normalization modes are currently implemented.

\paragraph{Value normalization.}
\begin{align}
  \psi_{\mathrm{norm}}(\mathbf{r})
  = \frac{\psi(\mathbf{r})}
         {\max_{\mathbf{r}} |\psi(\mathbf{r})| + \varepsilon}.
\end{align}

\paragraph{Reciprocal-space normalization.}
\begin{align}
  \psi_{\mathrm{norm}}(\mathbf{r})
  = \frac{\psi(\mathbf{r})}
         {\left(\sum_{\mathbf{k}} |\Psi(\mathbf{k})|^2\right)^{1/2} + \varepsilon}.
\end{align}

\subsection{PlaneWaveBuilder}
The plane-wave builder emits a constant incident field
\begin{align}
  \psi_0(\mathbf{r}) = A,
\end{align}
where $A \in \mathbb{C}$ is a user-selected complex amplitude.

\subsection{ProbeWavesBuilder}
The probe builder wraps the probe-forming optics and returns
\begin{align}
  \psi_{\mathrm{probe}}(\mathbf{r})
  = \mathcal{F}^{-1}\left\{ H(\mathbf{k}) \, A(\mathbf{k}) \right\},
\end{align}
where $A(\mathbf{k})$ is an optional reciprocal-space amplitude mask and
$H(\mathbf{k})$ is the transfer function defined below.

\section{Optics}
\subsection{CTF}
The contrast transfer function is modeled as
\begin{align}
  H(\mathbf{k})
  = P(\mathbf{k})
    \exp\!\left[- i \chi(\mathbf{k})\right]
    E_t(\mathbf{k})
    E_s(\mathbf{k}),
\end{align}
where $P$ is the aperture, $\chi$ is the aberration phase, $E_t$ is the temporal
envelope, and $E_s$ is the spatial envelope.  In the implementation, $E_t$ and $E_s$
are included only when the corresponding spread parameters are non-zero.

\paragraph{Relation to physical lens and aperture fields.}
The implemented CTF is a reduced effective optics model rather than a direct
electromagnetic field solve inside the microscope column.  The intended bridge to
hardware is
\begin{align}
  \bigl(\mathbf{E}(\mathbf{r}), \mathbf{B}(\mathbf{r}), \text{aperture geometry}\bigr)
  \;\Longrightarrow\;
  \text{optical path / eikonal phase}
  \;\Longrightarrow\;
  \chi(\mathbf{k}),\,P(\mathbf{k}).
\end{align}
At the Lorentz-force level, an electron ray or wave packet is driven by
\begin{align}
  \frac{d\mathbf{p}}{dt}
  =
  -e\left(
    \mathbf{E} + \mathbf{v}\times\mathbf{B}
  \right),
\end{align}
and, after the usual paraxial and reference-trajectory reduction, the accumulated
optical-path phase may be written schematically as
\begin{align}
  \Phi(\boldsymbol{\rho})
  \approx
  \Phi_0
  +
  \sum_j a_j\,f_j(\boldsymbol{\rho}),
\end{align}
where $\boldsymbol{\rho}$ is a pupil-plane coordinate and the coefficients $a_j$ are
low-order summaries of the underlying lens fields.  In reciprocal-space TEM notation,
this reduced phase is then represented by the aberration function
\begin{align}
  \chi(\mathbf{k}) = \Phi\!\left(\boldsymbol{\rho}(\mathbf{k})\right) - \Phi_0,
\end{align}
which is finally expanded in the aberration basis used by the simulator.  A real
aperture enters the same reduced model as the pupil support or amplitude mask
\begin{align}
  P(\mathbf{k}) \in [0,1],
\end{align}
with hard, soft, annular, segmented, or otherwise parameterized support depending on
the experiment branch.  So, in the current TorchTEM implementation, the aberration
coefficients and aperture masks should be interpreted as compressed observables of the
physical lens and aperture fields, not as a first-principles field map of the full
column.

Let $\alpha = 10^3 \lambda |\mathbf{k}|$ be the scattering angle in mrad and $\varphi$
the polar angle of $\mathbf{k}$.  The aberration phase polynomial is represented as
\begin{align}
  \chi(\alpha,\varphi)
  = \frac{2\pi}{\lambda}\sum_{n,m} \frac{\alpha^{n+1}}{n+1}
  C_{nm}\cos\!\left[m(\varphi-\phi_{nm})\right],
\end{align}
with explicit implemented terms up to fifth order.

The hard or soft aperture is
\begin{align}
  P(\alpha) =
  \begin{cases}
    \mathbf{1}\{\alpha \le \alpha_c\}, & \text{hard cutoff,} \\
    \sigma\!\left(\dfrac{\alpha_c-\alpha}{s}\right), & \text{soft edge,}
  \end{cases}
\end{align}
where $\alpha_c$ is the semiangle cutoff, $s$ is the soft-edge width, and $\sigma$ is
the logistic function.

The temporal envelope is
\begin{align}
  E_t(\alpha)
  = \exp\!\left[
    -\left(
      \frac{\pi}{2\lambda}\Delta f \, \alpha^2
    \right)^2
  \right],
\end{align}
where $\Delta f$ is the focal spread.

\section{Multislice Propagation}
\subsection{Transmission}
For slice $j$, the transmission function is
\begin{align}
  T_j(\mathbf{r}) = \exp\!\left[i \sigma V_j(\mathbf{r})\right].
\end{align}

\subsection{Fresnel propagation}
Between slices, the current implementation uses an exact Fourier-space propagator
\begin{align}
  \Psi_{j+1}(\mathbf{k})
  = G(\mathbf{k}; \Delta z)\,\Psi'_j(\mathbf{k}),
\end{align}
where $\Delta z$ is the slice thickness and
\begin{align}
  \Psi'_j(\mathbf{k}) = \mathcal{F}\left\{ T_j(\mathbf{r})\psi_j(\mathbf{r}) \right\}.
\end{align}
The propagated wave in real space is
\begin{align}
  \psi_{j+1}(\mathbf{r})
  = \mathcal{F}^{-1}\left\{
    G(\mathbf{k}; \Delta z)\,
    \mathcal{F}\left[T_j(\mathbf{r})\psi_j(\mathbf{r})\right]
  \right\}.
\end{align}

The implemented exact propagator distinguishes propagating and evanescent components via
$x = \lambda^2 |\mathbf{k}|^2$:
\begin{align}
  G(\mathbf{k};\Delta z) =
  \begin{cases}
    \exp\!\left[
      \dfrac{2\pi \Delta z}{\lambda}
      \left(
        -\dfrac{x}{\sqrt{1-x}+1}
      \right)
    \right], & x \le 1, \\
    \exp\!\left[
      \dfrac{2\pi \Delta z}{\lambda}
      \left(
        i\sqrt{x-1}-1
      \right)
    \right], & x > 1.
  \end{cases}
\end{align}

\subsection{Multislice recursion}
Starting from incident wave $\psi_0(\mathbf{r})$, the full multislice recursion is
\begin{align}
  \psi_{j+\frac{1}{2}}(\mathbf{r}) &= T_j(\mathbf{r})\psi_j(\mathbf{r}), \\
  \psi_{j+1}(\mathbf{r}) &= \mathcal{F}^{-1}\left\{
    G(\mathbf{k};\Delta z)\,
    \mathcal{F}\left[\psi_{j+\frac{1}{2}}(\mathbf{r})\right]
  \right\}.
\end{align}

\section{Wave--Optics Interaction}
Applying a CTF to a \texttt{Waves} object is defined as:
\begin{align}
  \Psi_{\mathrm{out}}(\mathbf{k})
  = H(\mathbf{k})\,\Psi_{\mathrm{in}}(\mathbf{k}),
\end{align}
where $\Psi_{\mathrm{in}}$ is either the stored reciprocal-space wave, or the Fourier
transform of a stored real-space wave.  If the input container is in real space, the
output is returned to real space after multiplication:
\begin{align}
  \psi_{\mathrm{out}}(\mathbf{r})
  = \mathcal{F}^{-1}\left\{
    H(\mathbf{k}) \mathcal{F}\{\psi_{\mathrm{in}}(\mathbf{r})\}
  \right\}.
\end{align}

\section{Specimen and Potential Layers}
\subsection{GaussianAtomProjection}
The current smooth surrogate projected-potential layer stores atom positions
$\mathbf{r}_n = (y_n, x_n)$, amplitudes $a_n$, and widths $\sigma_n$ and produces
\begin{align}
  V(\mathbf{r})
  = \sum_{n=1}^{N_{\mathrm{atom}}}
    a_n
    \exp\!\left(
      -\frac{\|\mathbf{r}-\mathbf{r}_n\|^2}{2\sigma_n^2}
    \right).
\end{align}
This layer is used as a trainable, differentiable surrogate while more faithful
parametrized atomic potentials are also available elsewhere in the implementation.

\subsection{Projected potential slices}
When a single projected potential $V(\mathbf{r})$ is reused across a multislice stack of
$N_s$ equal slices, the current \texttt{TEMModel} path distributes it as
\begin{align}
  V_j(\mathbf{r}) = \frac{1}{N_s}V(\mathbf{r}),
  \qquad j=1,\dots,N_s.
\end{align}
If the potential module already returns a rank-3 tensor of slice-resolved potentials,
that tensor is interpreted directly as
\begin{align}
  \{V_j(\mathbf{r})\}_{j=1}^{N_s}.
\end{align}

\subsection{IAM slice builder}
The current IAM potential builder stores atoms
\begin{align}
  \mathbf{R}_n = (y_n, x_n, z_n),
\end{align}
chemical symbols $s_n$, and optional amplitudes $a_n$, and returns a slice stack
\begin{align}
  \{V_j(\mathbf{r})\}_{j=1}^{N_s},
  \qquad
  [z_j^-, z_j^+) \subset [0, L_z).
\end{align}

For infinite projected-potential mode, each atom is assigned to exactly one slice by
\begin{align}
  m_j(n) = \mathbf{1}\{z_n \in [z_j^-, z_j^+)\},
\end{align}
and the slice potential is assembled as
\begin{align}
  V_j(\mathbf{r})
  = \sum_{n=1}^{N_{\mathrm{atom}}}
    m_j(n)\,a_n\,\Phi_{s_n}^{\infty}(\mathbf{r} - \mathbf{R}_{n,xy}),
\end{align}
where $\Phi_{s_n}^{\infty}$ denotes the chosen infinite projected atomic potential,
either from the direct parametrized projected-potential formulas or from the
projected-scattering-factor integrator path.

For finite projected mode, each slice accumulates the finite $z$-window contribution
of every atom,
\begin{align}
  V_j(\mathbf{r})
  = \sum_{n=1}^{N_{\mathrm{atom}}}
    a_n\,\Phi_{s_n}^{[z_j^-, z_j^+)}(\mathbf{r} - \mathbf{R}_{n,xy}, z_n),
\end{align}
where $\Phi_{s_n}^{[z_j^-, z_j^+)}$ is evaluated either by the explicit finite Peng
formula or by the Gaussian-plus-correction projection-integral route used for the
current Lobato/Kirkland finite builder support.

\section{Scan Layers}
\subsection{Fourier-shift scan positioning}
Given a reference wave $\psi(\mathbf{r})$ and scan shift
$\Delta \mathbf{r}_m = (\Delta y_m, \Delta x_m)$, the implementation applies the Fourier
shift theorem:
\begin{align}
  \psi_m(\mathbf{r})
  = \mathcal{F}^{-1}\left\{
    \exp\!\left[-2\pi i
      \left(
        \Delta y_m k_y + \Delta x_m k_x
      \right)
    \right]
    \mathcal{F}\{\psi(\mathbf{r})\}
  \right\}.
\end{align}

\subsection{GridScan}
For scan start point $\mathbf{r}_{\min} = (y_{\min},x_{\min})$, end point
$\mathbf{r}_{\max} = (y_{\max},x_{\max})$, and raster shape
$(N_y^{\mathrm{scan}},N_x^{\mathrm{scan}})$, the scan positions are
\begin{align}
  y_i &= y_{\min} + \frac{i}{N_y^{\mathrm{scan}}-1}(y_{\max}-y_{\min}), \\
  x_j &= x_{\min} + \frac{j}{N_x^{\mathrm{scan}}-1}(x_{\max}-x_{\min}),
\end{align}
with lexicographically flattened position list
\begin{align}
  \mathcal{S}
  = \left\{ (y_i,x_j) \right\}_{i,j}.
\end{align}

\subsection{CustomScan}
The custom scan layer is a direct position map:
\begin{align}
  \mathcal{S}
  = \left\{ (y_m,x_m) \right\}_{m=1}^{M},
\end{align}
where the full list is supplied by the caller.

\section{Detection Layers}
\subsection{DiffractionIntensity}
The diffraction intensity operator computes
\begin{align}
  I(\mathbf{k}) = |\mathcal{F}\{\psi(\mathbf{r})\}|^2.
\end{align}

\subsection{AnnularDetector}
Let the angle components in mrad be
\begin{align}
  \alpha_y(\mathbf{k}) &= 10^3 \lambda k_y, \\
  \alpha_x(\mathbf{k}) &= 10^3 \lambda k_x,
\end{align}
and let $\Delta\boldsymbol{\alpha}=(\Delta\alpha_y,\Delta\alpha_x)$ be an optional
detector-center offset.  The shifted radial coordinate is
\begin{align}
  \alpha_\Delta(\mathbf{k})
  =
  \sqrt{
    \left(\alpha_y(\mathbf{k})-\Delta\alpha_y\right)^2
    +
    \left(\alpha_x(\mathbf{k})-\Delta\alpha_x\right)^2
  }.
\end{align}
For inner and outer detector radii $\alpha_{\mathrm{in}}$ and
$\alpha_{\mathrm{out}}$, the annular mask is
\begin{align}
  M_{\mathrm{ann}}(\mathbf{k})
  = \mathbf{1}\{
    \alpha_{\mathrm{in}} \le \alpha_\Delta(\mathbf{k}) < \alpha_{\mathrm{out}}
  \},
\end{align}
and the detector output is
\begin{align}
  D_{\mathrm{ann}}
  = \sum_{\mathbf{k}} I(\mathbf{k}) M_{\mathrm{ann}}(\mathbf{k}).
\end{align}
The annular detector-region helper returns the binary mask
$M_{\mathrm{ann}}(\mathbf{k})$ itself for geometry inspection.

\subsection{PixelatedDetector}
The pixelated detector may either record diffraction intensity
\begin{align}
  D_{\mathrm{pix}}(\mathbf{k}) = I(\mathbf{k}).
\end{align}
or real-space wave intensity
\begin{align}
  D_{\mathrm{pix,real}}(\mathbf{r}) = |\psi(\mathbf{r})|^2.
\end{align}
If a target output grid $\mathbf{N}'$ is requested, the current Torch-native
implementation resamples the intensity image to that shape using the shared FFT-based
interpolation helper.
If a maximum detected scattering angle $\alpha_{\max}$ is configured in reciprocal
mode, the detector first applies the mask
\begin{align}
  M_{\max}(\mathbf{k}) = \mathbf{1}\{\alpha(\mathbf{k}) \le \alpha_{\max}\},
\end{align}
and returns
\begin{align}
  D_{\mathrm{pix},\max}(\mathbf{k}) = I(\mathbf{k}) M_{\max}(\mathbf{k}).
\end{align}
The corresponding detector-region helper returns $M_{\max}(\mathbf{k})$ directly; in
real-space mode it returns the all-ones mask because the full image plane is accepted.

\subsection{SegmentedDetector}
Let the valid angular interval be
\begin{align}
  \alpha_{\mathrm{in}} \le \alpha_\Delta(\mathbf{k}) < \alpha_{\mathrm{out}},
\end{align}
and let the shifted polar angle be
\begin{align}
  \varphi_\Delta(\mathbf{k})
  =
  \mathrm{mod}\!\left(
    \mathrm{atan2}\!\left(
      \alpha_x(\mathbf{k})-\Delta\alpha_x,
      \alpha_y(\mathbf{k})-\Delta\alpha_y
    \right)
    - \varphi_0,
    2\pi
  \right),
\end{align}
where $\varphi_0$ is the detector rotation.  The radial and azimuthal bin indices are
\begin{align}
  b_r(\mathbf{k}) &=
  \left\lfloor
    N_r \frac{\alpha_\Delta(\mathbf{k})-\alpha_{\mathrm{in}}}
               {\alpha_{\mathrm{out}}-\alpha_{\mathrm{in}}}
  \right\rfloor, \\
  b_\varphi(\mathbf{k}) &=
  \left\lfloor
    N_\varphi \frac{\varphi_\Delta(\mathbf{k})}{2\pi}
  \right\rfloor.
\end{align}
The segment index is
\begin{align}
  b(\mathbf{k}) = b_r(\mathbf{k}) N_\varphi + b_\varphi(\mathbf{k}),
\end{align}
and the detector output accumulates intensity over each segment:
\begin{align}
  D_{r,\varphi}
  = \sum_{\mathbf{k} \in \Omega_{r,\varphi}} I(\mathbf{k}).
\end{align}
The segmented detector-region helper returns the family of masks
$\{\mathbf{1}_{\Omega_{r,\varphi}}(\mathbf{k})\}_{r,\varphi}$, represented with
NaN outside each region so per-region plots remain visually separated.

\section{Coherence Averaging}
The current source-offset ensemble path evaluates multiple coherent members
\begin{align}
  \{\psi^{(m)}(\mathbf{r})\}_{m=1}^{M}
\end{align}
and averages their measured intensities:
\begin{align}
  \bar{D}
  = \frac{1}{M}\sum_{m=1}^{M} D\!\left[\psi^{(m)}\right].
\end{align}
This corresponds to an incoherent mixture approximation for the selected source
ensemble.

\section{Assembly and Composition Layers}
\subsection{LayerStack}
The stack assembly API represents the full simulator as an ordered composition
\begin{align}
  y = (f_L \circ f_{L-1} \circ \cdots \circ f_1 \circ s)(\varnothing),
\end{align}
where $s$ is a source layer with no wave input and each $f_\ell$ is a unary map on
either a wave or a derived measurement tensor.  The implementation also allows the
intermediate states
\begin{align}
  z_0 &= s(\varnothing), \\
  z_\ell &= f_\ell(z_{\ell-1}), \qquad \ell=1,\dots,L,
\end{align}
to be returned explicitly for debugging, supervision, or auxiliary losses.

\subsection{PotentialSliceInteraction}
The potential-slice interaction layer binds a potential-producing module to a
multislice propagator:
\begin{align}
  \psi_{\mathrm{out}}
  = \mathcal{M}\!\left(
    \psi_{\mathrm{in}},
    \{V_j\}_{j=1}^{N_s}
  \right),
\end{align}
where $\mathcal{M}$ denotes the multislice recursion and the potential module produces
either a slice stack $\{V_j\}_{j=1}^{N_s}$ directly or a single projected potential
$V(\mathbf{r})$ that is distributed uniformly as
\begin{align}
  V_j(\mathbf{r}) = \frac{1}{N_s}V(\mathbf{r}).
\end{align}

\subsection{ScanLayer}
The scan layer converts a reference wave into a batch of shifted waves over a position
set $\mathcal{S} = \{\Delta \mathbf{r}_m\}_{m=1}^{M}$:
\begin{align}
  \psi^{(m)}(\mathbf{r})
  = \mathcal{T}_{\Delta \mathbf{r}_m}\psi(\mathbf{r}),
  \qquad m=1,\dots,M,
\end{align}
where $\mathcal{T}_{\Delta \mathbf{r}_m}$ is the implemented Fourier-shift operator
from the scan section above.  The output is therefore a stacked tensor
\begin{align}
  \Psi_{\mathrm{scan}}
  = \{\psi^{(m)}\}_{m=1}^{M}.
\end{align}

\subsection{DetectorReadout and BatchDetectorReadout}
For a detector functional $\mathcal{D}$, the unary detector readout is simply
\begin{align}
  y = \mathcal{D}[\psi].
\end{align}
For a leading scan or ensemble batch, the batch detector path applies the same
detector independently:
\begin{align}
  y_m = \mathcal{D}[\psi^{(m)}],
  \qquad m=1,\dots,M.
\end{align}
This keeps the whole assembled simulator differentiable as a map from source and
specimen parameters to either scalar signals, image stacks, or diffraction data.

\subsection{Structured experiment assembly}
The experiment-builder layer defines a parameter object
\begin{align}
  \Theta
  = \left(
    \Theta_{\mathrm{src}},
    \Theta_{\mathrm{ms}},
    \Theta_{\mathrm{pot}},
    \Theta_{\mathrm{scan}},
    \Theta_{\mathrm{det}},
    \Theta_{\mathrm{coh}}
  \right),
\end{align}
where omitted blocks are allowed for modes that do not require them.  The builder then
instantiates a composed simulator
\begin{align}
  \mathcal{E}_{\Theta}
  = \mathcal{R}_{\Theta_{\mathrm{coh}},\Theta_{\mathrm{det}}}
    \circ
    \mathcal{S}_{\Theta_{\mathrm{scan}}}
    \circ
    \mathcal{M}_{\Theta_{\mathrm{ms}},\Theta_{\mathrm{pot}}}
    \circ
    \mathcal{Q}_{\Theta_{\mathrm{src}}},
\end{align}
with the conventions that $\mathcal{S}$ becomes the identity if no scan is present,
and $\mathcal{R}$ reduces to the identity, detector-only readout, or coherence-aware
readout depending on the selected experiment mode.

For a plane-wave source configuration,
\begin{align}
  \mathcal{Q}_{\Theta_{\mathrm{src}}}(\varnothing)
  = A \mathbf{1},
\end{align}
whereas for a probe-forming source configuration,
\begin{align}
  \mathcal{Q}_{\Theta_{\mathrm{src}}}(\varnothing)
  = \mathcal{F}^{-1}\left\{
    H_{\Theta_{\mathrm{src}}}(\mathbf{k})
  \right\}.
\end{align}

The resulting experiment map therefore provides a differentiable parameter-to-output
operator
\begin{align}
  y = \mathcal{E}_{\Theta}(\varnothing),
\end{align}
which is the object intended for downstream gradient-based calibration, inverse
problems, and external training pipelines.

\subsection{Multi-detector experiment branches}
For realistic STEM/TEM experiments the detector block may itself be a family of
simultaneous readouts
\begin{align}
  \mathcal{D}
  = \left\{
    \mathcal{D}_q
  \right\}_{q \in \mathcal{Q}},
\end{align}
where $\mathcal{Q}$ is a set of user-defined detector names such as annular,
segmented, or pixelated channels.  The experiment readout then returns a named output
map
\begin{align}
  y_q = \mathcal{D}_q[\psi],
  \qquad q \in \mathcal{Q},
\end{align}
or, for scanned experiments,
\begin{align}
  y_q^{(m)} = \mathcal{D}_q[\psi^{(m)}],
  \qquad m=1,\dots,M.
\end{align}
With coherence averaging enabled in the ensemble-averaging mode, each named branch becomes
\begin{align}
  y_q^{(m)}
  = \sum_{\ell=1}^{L} w_\ell \,
    \mathcal{D}_q\!\left[\psi_\ell^{(m)}\right],
\end{align}
with normalized coherence weights $w_\ell$.  This is the semantics used by the
Torch-native multi-detector experiment builder and plotting surface.

For diagonal mixed coherence, the experiment source is replaced by coherent modes
$\psi_\ell^{(0)}$ generated by the diagonal mixed-coherence factor decomposition,
the elastic stack propagates each mode independently,
\begin{align}
  \psi_\ell^{(m)}
  = \mathcal{M}_{\Theta_{\mathrm{ms}},\Theta_{\mathrm{pot}}}
    \mathcal{S}_{\mathbf{r}^{(m)}} \psi_\ell^{(0)},
\end{align}
and the detector outputs are reduced incoherently after propagation:
\begin{align}
  y_q^{(m)}
  = \sum_{\ell=1}^{L}
    \mathcal{D}_q\!\left[\psi_\ell^{(m)}\right].
\end{align}
If no detector is configured, the experiment returns the incoherent real-space intensity
\begin{align}
  I^{(m)}(\mathbf{r})
  = \sum_{\ell=1}^{L} \left|\psi_\ell^{(m)}(\mathbf{r})\right|^2.
\end{align}

\subsection{Mode presets}
The preset experiment surfaces are typed restrictions of the general experiment map
$\mathcal{E}_{\Theta}$.

\paragraph{HRTEM preset.}
The HRTEM preset omits scanning and uses
\begin{align}
  y = \mathcal{D}\!\left[
    \mathcal{M}_{\Theta_{\mathrm{ms}},\Theta_{\mathrm{pot}}}
    \mathcal{Q}_{\Theta_{\mathrm{src}}}(\varnothing)
  \right],
\end{align}
where $\mathcal{D}$ is typically a pixelated image-plane or diffraction-plane readout.

\paragraph{STEM preset.}
The STEM preset requires a probe-forming source and scan path:
\begin{align}
  y^{(m)} = \mathcal{D}\!\left[
    \mathcal{M}_{\Theta_{\mathrm{ms}},\Theta_{\mathrm{pot}}}
    \mathcal{T}_{\Delta \mathbf{r}_m}
    \mathcal{Q}_{\Theta_{\mathrm{src}}}(\varnothing)
  \right].
\end{align}

\paragraph{Diffraction preset.}
The diffraction preset uses a pixelated detector branch by default:
\begin{align}
  y^{(m)}(\mathbf{k}) = \left|
    \mathcal{F}\left\{
      \mathcal{M}_{\Theta_{\mathrm{ms}},\Theta_{\mathrm{pot}}}
      \mathcal{T}_{\Delta \mathbf{r}_m}
      \mathcal{Q}_{\Theta_{\mathrm{src}}}(\varnothing)
    \right\}
  \right|^2,
\end{align}
with the scan operator removed when no scan path is configured.

\subsection{Simulation-result wrapper}
The simulator-facing run helper does not change the computed output; it packages the
same final tensor or named detector map together with metadata and optional
intermediate states.  If the layer stack has states
\begin{align}
  \{z_\ell\}_{\ell=0}^{L},
\end{align}
then the wrapped result stores
\begin{align}
  y &= z_L, \\
  \mathcal{I} &= \{z_\ell\}_{\ell=0}^{L},
\end{align}
when intermediate capture is requested, and otherwise stores only $y$ together with
metadata such as experiment mode, scan/coherence flags, and detector names.  Plotting
is defined as a visualization map
\begin{align}
  \mathcal{P}: y \mapsto \text{matplotlib panels},
\end{align}
where each detector branch or scan-indexed slice is rendered as a labeled panel.

\subsection{Post-detector measurement pipelines}
The experiment surface may append a post-detector pipeline
\begin{align}
  \mathcal{R} = \mathcal{R}_P \circ \cdots \circ \mathcal{R}_2 \circ \mathcal{R}_1
\end{align}
after detector readout.  The resulting experiment output becomes
\begin{align}
  \tilde{y} = \mathcal{R}(y),
\end{align}
or, for named detector branches,
\begin{align}
  \tilde{y}_q = \mathcal{R}_q(y_q), \qquad q \in \mathcal{Q},
\end{align}
where each branch may have its own configured sequence.

Implemented examples include:
\begin{align}
  \mathcal{R}_{\mathrm{blur}}(I) &= G_\sigma * I, \\
  \mathcal{R}_{\mathrm{MTF}}(I) &= \mathcal{F}^{-1}\left\{
    \sqrt{K_{\mathrm{MTF}}(\mathbf{k})}\,\mathcal{F}\{I\}
  \right\}, \\
  \mathcal{R}_{\mathrm{Poisson}}(I) &\sim \mathrm{Poisson}(d\,I), \\
  \mathcal{R}_{\mathrm{COM}}(I) &= \left(
    \frac{\sum_{\mathbf{k}} k_y I(\mathbf{k})}{\sum_{\mathbf{k}} I(\mathbf{k})},
    \frac{\sum_{\mathbf{k}} k_x I(\mathbf{k})}{\sum_{\mathbf{k}} I(\mathbf{k})}
  \right).
\end{align}
This is the semantics used by the Torch-native experiment postprocessing layer.

\subsection{Propagation-backend selection}
The experiment multislice configuration may choose between distinct propagation
operators while preserving the same higher-level assembly semantics.  Denote the chosen
propagation backend by $\mathcal{M}^{(b)}$, where
\begin{align}
  b \in \{\text{fourier}, \text{realspace}\}.
\end{align}
The experiment map then uses
\begin{align}
  y = \mathcal{D}\!\left[
    \mathcal{M}^{(b)}_{\Theta_{\mathrm{ms}},\Theta_{\mathrm{pot}}}
    \mathcal{Q}_{\Theta_{\mathrm{src}}}(\varnothing)
  \right].
\end{align}
For $b=\text{fourier}$ this is the reciprocal-space Fresnel multislice recursion
described above.  For $b=\text{realspace}$ it is the truncated exponential-series
real-space propagator already implemented in the low-level stack.  The experiment layer
therefore selects among low-level propagators without changing the surrounding source,
scan, detector, or measurement-pipeline structure.

For $b=\text{prism}$, the scanned probe experiment uses an S-matrix reduction path:
\begin{align}
  S_p(\mathbf{r})
  &=
  \mathcal{M}^{(\text{fourier})}_{\Theta_{\mathrm{ms}},\Theta_{\mathrm{pot}}}
  \phi_p(\mathbf{r}), \\
  \psi^{(m)}(\mathbf{r})
  &= \sum_{p=1}^{P}
    c_p \,
    \exp\!\left[
      -2\pi i \,
      \mathbf{k}_p \cdot \Delta \mathbf{r}_m
    \right]
    S_p(\mathbf{r}),
\end{align}
where $\phi_p$ are parent plane waves, $\mathbf{k}_p$ are the selected parent wave
vectors, and $c_p$ are the probe coefficients sampled from the CTF grid.  For a scan
configuration, the phase factor produces one reduced wave per scan position
$\Delta \mathbf{r}_m$; without a scan configuration, the current experiment layer
reduces the S-matrix once at the origin. This is the current experiment-level PRISM
path exposed by the Torch-native simulator surface.

For $b=\text{bloch}$, the experiment currently exposes a diffraction-style path driven
by structure factors:
\begin{align}
  \mathbf{c}(t) &= S_{\mathrm{Bloch}}(t)\,\mathbf{c}_0, \\
  I_j &= |c_j(t)|^2,
\end{align}
where $\mathbf{c}_0$ is the incident plane-wave coefficient vector and
$S_{\mathrm{Bloch}}$ is the Bloch scattering matrix from the low-level model.  The
experiment layer rasterizes $\{I_j\}$ onto the configured reciprocal-space image grid
at the corresponding projected reciprocal-lattice coordinates.  Reciprocal-intensity
detector readout may then be applied directly to this rasterized diffraction image for
the current pixelated-reciprocal, annular, flexible-annular, and segmented detector
families.  Alternatively, the complex coefficients $c_j(t)$ themselves may be
rasterized onto the reciprocal grid for reciprocal-wave output. This is intentionally a
constrained diffraction-oriented surface rather than full Bloch feature parity.

\subsection{Inelastic branch}
The experiment layer may append a constrained inelastic branch after elastic
propagation, and this branch may also be followed by the current coherence-averaging
readout path:
\begin{align}
  \psi_{\mathrm{el}}
  = \mathcal{M}_{\Theta_{\mathrm{ms}},\Theta_{\mathrm{pot}}}
    \mathcal{Q}_{\Theta_{\mathrm{src}}}(\varnothing),
\end{align}
followed by an inelastic operator $\mathcal{T}_{\mathrm{inel}}$.

For channel-producing transition potentials,
\begin{align}
  \psi_c^{\mathrm{inel}}(\mathbf{r})
  = T_c(\mathbf{r}) \, \psi_{\mathrm{el}}(\mathbf{r}),
\end{align}
and the experiment output is either the stack $\{\psi_c^{\mathrm{inel}}\}_c$ or, when a
detector branch is configured, the recursively reduced detector outputs
\begin{align}
  y_c = \mathcal{D}\!\left[\psi_c^{\mathrm{inel}}\right].
\end{align}

For plasmon weighted-intensity mode, the branch evaluates shifted inelastic members
\begin{align}
  \psi_\ell^{\mathrm{pl}}(\mathbf{r})
  = \mathcal{S}_{\Delta \mathbf{r}_\ell}\psi_{\mathrm{el}}(\mathbf{r}),
\end{align}
and returns
\begin{align}
  I_{\mathrm{pl}}(\mathbf{r})
  = \sum_{\ell} w_\ell \left|\psi_\ell^{\mathrm{pl}}(\mathbf{r})\right|^2.
\end{align}
If a coherence model is configured, detector or intensity outputs are reduced over
coherent members after the inelastic branch using the same ensemble-weighted or
diagonal-MCF incoherent reduction as the elastic experiment path. This is the current
constrained inelastic experiment surface exposed by TorchTEM.

\subsection{Magnetic branch}
The experiment layer may also append a projected magnetic interaction after elastic
propagation:
\begin{align}
  \psi_{\mathrm{mag}}
  = \mathcal{P}_A\!\left(
    A_{xy},
    \psi_{\mathrm{el}}
  \right),
\end{align}
where $A_{xy}$ is a projected vector potential from the magnetic builder and
$\mathcal{P}_A$ is the projected Pauli interaction currently implemented as
\begin{align}
  \psi_{\mathrm{mag}}
  = \psi_{\mathrm{el}} + \eta \, A_{xy}\cdot\nabla_{xy}\psi_{\mathrm{el}},
\end{align}
with configurable complex prefactor $\eta$.  The resulting magnetic wave may then feed
the usual detector and postprocessing branches.  This is the current constrained
magnetic experiment surface exposed by TorchTEM.

\subsection{Image-forming detector branch}
The Torch-native detector family now also includes a real-space image-forming branch.
Given an exit wave $\psi_{\mathrm{exit}}(\mathbf{r})$, the detector may optionally
apply an objective-lens transfer function
\begin{align}
  H(\mathbf{k}) = A(\mathbf{k})\,\exp\!\bigl[-i\chi(\mathbf{k})\bigr],
\end{align}
where $A$ is the aperture and $\chi$ is the aberration phase polynomial from the CTF
layer already defined above.  The image-plane wave is then
\begin{align}
  \psi_{\mathrm{img}}(\mathbf{r})
  =
  \mathcal{F}^{-1}
  \left\{
    H(\mathbf{k})\,
    \mathcal{F}\{\psi_{\mathrm{exit}}(\mathbf{r})\}
  \right\}.
\end{align}
The detector may return several real-space channels:
\begin{align}
  I(\mathbf{r}) &= |\psi_{\mathrm{img}}(\mathbf{r})|^2, \\
  \varphi(\mathbf{r}) &= \arg \psi_{\mathrm{img}}(\mathbf{r}), \\
  R(\mathbf{r}) &= \Re \psi_{\mathrm{img}}(\mathbf{r}), \\
  J(\mathbf{r}) &= \Im \psi_{\mathrm{img}}(\mathbf{r}),
\end{align}
or the complex image-plane wave itself.  In experiment-builder notation, this makes it
possible to express HRTEM-style image formation as a detector branch
\begin{align}
  y = \mathcal{D}_{\mathrm{img}}
  \left[
    \mathcal{M}_{\Theta_{\mathrm{ms}},\Theta_{\mathrm{pot}}}
    \mathcal{Q}_{\Theta_{\mathrm{src}}}(\varnothing)
  \right]
\end{align}
instead of applying an imaging CTF manually outside the assembled simulator graph.

\subsection{Direct waves detector branch}
The detector family also includes a direct wave readout branch for optimization or
reconstruction pipelines that require complex fields instead of already reduced image
intensities.  Given an exit wave $\psi_{\mathrm{exit}}(\mathbf{r})$, the detector may
either return the real-space wave directly,
\begin{align}
  y_{\mathrm{exit}}(\mathbf{r}) = \psi_{\mathrm{exit}}(\mathbf{r}),
\end{align}
or the reciprocal-space wave
\begin{align}
  y_{\mathrm{rec}}(\mathbf{k}) = \mathcal{F}\{\psi_{\mathrm{exit}}(\mathbf{r})\}.
\end{align}
If a target detector grid $\mathbf{N}'=(N_y',N_x')$ is requested, the returned wave is
resampled by FFT cropping/interpolation:
\begin{align}
  \tilde{y} = \mathcal{I}_{\mathbf{N}'}[y],
\end{align}
where $\mathcal{I}_{\mathbf{N}'}$ denotes the Torch-native Fourier interpolation
operator already implemented in the support layer.  This yields a detector branch
\begin{align}
  y = \mathcal{D}_{\mathrm{waves}}
  \left[
    \mathcal{M}_{\Theta_{\mathrm{ms}},\Theta_{\mathrm{pot}}}
    \mathcal{Q}_{\Theta_{\mathrm{src}}}(\varnothing)
  \right]
\end{align}
that preserves the complex field for downstream differentiable pipelines.

\subsection{Flexible annular detector branch}
The detector family also includes a radial-binning annular detector that preserves
annular integration flexibility after simulation.  Let the diffraction intensity be
\begin{align}
  I(\mathbf{k}) = \left|\mathcal{F}\{\psi(\mathbf{r})\}\right|^2,
\end{align}
and define the shifted scattering-angle magnitude
\begin{align}
  \alpha_\Delta(\mathbf{k}) =
  \sqrt{
    \left(\alpha_y(\mathbf{k})-\Delta\alpha_y\right)^2
    +
    \left(\alpha_x(\mathbf{k})-\Delta\alpha_x\right)^2
  },
\end{align}
in mrad.  For inner angle $\alpha_{\mathrm{in}}$, optional outer angle
$\alpha_{\mathrm{out}}$, and radial step size $\Delta \alpha$, the detector defines
radial bin indices
\begin{align}
  b(\mathbf{k}) =
  \left\lfloor
    \frac{\alpha_\Delta(\mathbf{k}) - \alpha_{\mathrm{in}}}{\Delta \alpha}
  \right\rfloor,
\end{align}
restricted to pixels satisfying
\begin{align}
  \alpha_{\mathrm{in}} \le \alpha_\Delta(\mathbf{k}) < \alpha_{\mathrm{out}}.
\end{align}
The detector output is then the vector of radial annular sums
\begin{align}
  y_n = \sum_{\mathbf{k}: b(\mathbf{k}) = n} I(\mathbf{k}),
\end{align}
for $n=0,\dots,N_r-1$, where
\begin{align}
  N_r = \left\lfloor
    \frac{\alpha_{\mathrm{out}} - \alpha_{\mathrm{in}}}{\Delta \alpha}
  \right\rfloor.
\end{align}
This provides a differentiable analogue of the original flexible annular detector:
the simulation produces a radial-binned detector signal once, and downstream annular
integrations may be chosen afterward without recomputing the forward wave propagation.
The Torch-native implementation now also supports off-axis detector-center shifts
$\Delta\boldsymbol{\alpha}$ in this branch.
The corresponding detector-region helper returns the integer-valued radial bin map
$b(\mathbf{k})$ with invalid pixels marked outside the accepted annular support.
An additional helper returns the family of per-bin masks
$\{\mathbf{1}_{b(\mathbf{k})=n}\}_n$, represented as separate 2D region planes with
NaN outside the active bin so radial detector regions can be visualized directly.
Because the detector output is the radial-bin vector $\{y_n\}$, a post-hoc annular
integration over limits $(\alpha_1,\alpha_2)$ is implemented as
\begin{align}
  D_{\mathrm{flex}}(\alpha_1,\alpha_2)
  = \sum_{n=n_1}^{n_2-1} y_n,
\end{align}
where
\begin{align}
  n_1 &= \left\lfloor
    \frac{\alpha_1-\alpha_{\mathrm{in}}}{\Delta \alpha}
  \right\rfloor, \\
  n_2 &= \left\lceil
    \frac{\alpha_2-\alpha_{\mathrm{in}}}{\Delta \alpha}
  \right\rceil.
\end{align}
For limits aligned with the radial bin boundaries, this matches the corresponding fixed
annular detector integral exactly up to floating-point arithmetic.


\subsection{Fixed-stack HRTEM control-vector branch}
TorchTEM also includes a fixed-topology HRTEM branch for inverse microscope tuning with
a constant specimen.  Let
\begin{align}
  \mathbf{u}\in\mathbb{R}^{12}
\end{align}
denote the control vector
\begin{align}
  \mathbf{u} =
  (&I_{c1}, I_{c2}, s_{c,x}, s_{c,y}, f_o, s_{o,x}, s_{o,y}, \gamma_{30},
    \gamma_{32,x}, \gamma_{32,y}, \gamma_{34,x}, \gamma_{34,y}),
\end{align}
where the entries represent two condenser excitations, condenser stigmator components,
objective focus, objective stigmator components, and simple corrector controls.
For batched inference or optimization, the implementation also accepts a tensor
\begin{align}
  U \in \mathbb{R}^{B_1\times\cdots\times B_m\times 12},
\end{align}
which is interpreted as a batch of control vectors with arbitrary leading batch shape.
The current fixed-stack surface also groups these entries into stage blocks
\begin{align}
  \mathbf{u}
  =
  \left(
    \mathbf{u}_{\mathrm{cond}},
    \mathbf{u}_{\mathrm{obj}},
    \mathbf{u}_{\mathrm{corr}}
  \right),
\end{align}
with
\begin{align}
  \mathbf{u}_{\mathrm{cond}} &\in \mathbb{R}^{4}, \\
  \mathbf{u}_{\mathrm{obj}} &\in \mathbb{R}^{3}, \\
  \mathbf{u}_{\mathrm{corr}} &\in \mathbb{R}^{5},
\end{align}
corresponding respectively to the condenser, objective, and corrector stage vectors.
In the implementation these stage blocks may be extracted or replaced independently
before running the full differentiable HRTEM forward model.
For inverse fitting, the current stack also supports an active subspace
\begin{align}
  \mathbf{u}_{\mathcal{A}} \in \mathbb{R}^{|\mathcal{A}|},
\end{align}
which replaces only a selected subset of control entries while the remaining entries
stay fixed at a baseline microscope state.  When bounds
\(
  \ell_i < u_i < h_i
\)
are supplied for an active coordinate, the trainable latent parameter
\(
  z_i \in \mathbb{R}
\)
is mapped to the physical control value by
\begin{align}
  u_i = \ell_i + (h_i - \ell_i)\,\sigma(z_i),
\end{align}
with \(\sigma\) the logistic sigmoid.  Unbounded active coordinates are represented
directly.  This keeps the differentiable fitting path inside a chosen physically
plausible parameter window without changing the fixed-stack forward model itself.

The current demo implementation decodes this vector into source-side and image-side
aberration coefficients through an explicit affine calibration map.  For the source
branch, the most important terms are
\begin{align}
  C_{10}^{(\mathrm{src})} &= -65 I_{c1} - 28 I_{c2}, \\
  a_{12,x}^{(\mathrm{src})} &= 24 s_{c,x} + 7 s_{c,y}, \\
  a_{12,y}^{(\mathrm{src})} &= 7 s_{c,x} - 24 s_{c,y}, \\
  C_{12}^{(\mathrm{src})} &=
  \sqrt{ \left(a_{12,x}^{(\mathrm{src})}\right)^2
       + \left(a_{12,y}^{(\mathrm{src})}\right)^2 }, \\
  \phi_{12}^{(\mathrm{src})} &=
  \tfrac{1}{2}\operatorname{atan2}
  \left(a_{12,y}^{(\mathrm{src})}, a_{12,x}^{(\mathrm{src})}\right), \\
  C_{21}^{(\mathrm{src})} &= 14\,(I_{c1}-I_{c2}).
\end{align}
For the image-forming branch, the current demo map uses
\begin{align}
  C_{10}^{(\mathrm{img})} &= 110 f_o + 6 (I_{c1}-I_{c2}), \\
  a_{12,x}^{(\mathrm{img})} &= 18 s_{o,x} - 900 \gamma_{32,x}, \\
  a_{12,y}^{(\mathrm{img})} &= 18 s_{o,y} - 900 \gamma_{32,y}, \\
  C_{12}^{(\mathrm{img})} &=
  \sqrt{ \left(a_{12,x}^{(\mathrm{img})}\right)^2
       + \left(a_{12,y}^{(\mathrm{img})}\right)^2 }, \\
  \phi_{12}^{(\mathrm{img})} &=
  \tfrac{1}{2}\operatorname{atan2}
  \left(a_{12,y}^{(\mathrm{img})}, a_{12,x}^{(\mathrm{img})}\right), \\
  C_{30}^{(\mathrm{img})} &= 1.2\times 10^4 - 8.5\times 10^3 \gamma_{30}, \\
  a_{32,x}^{(\mathrm{img})} &= -2200 \gamma_{32,x}, \\
  a_{32,y}^{(\mathrm{img})} &= -2200 \gamma_{32,y}, \\
  C_{32}^{(\mathrm{img})} &=
  \sqrt{ \left(a_{32,x}^{(\mathrm{img})}\right)^2
       + \left(a_{32,y}^{(\mathrm{img})}\right)^2 }, \\
  \phi_{32}^{(\mathrm{img})} &=
  \tfrac{1}{2}\operatorname{atan2}
  \left(a_{32,y}^{(\mathrm{img})}, a_{32,x}^{(\mathrm{img})}\right), \\
  a_{34,x}^{(\mathrm{img})} &= -1400 \gamma_{34,x}, \\
  a_{34,y}^{(\mathrm{img})} &= -1400 \gamma_{34,y}, \\
  C_{34}^{(\mathrm{img})} &=
  \sqrt{ \left(a_{34,x}^{(\mathrm{img})}\right)^2
       + \left(a_{34,y}^{(\mathrm{img})}\right)^2 }, \\
  \phi_{34}^{(\mathrm{img})} &=
  \tfrac{1}{4}\operatorname{atan2}
  \left(a_{34,y}^{(\mathrm{img})}, a_{34,x}^{(\mathrm{img})}\right).
\end{align}
All higher coefficients not listed above are currently set to zero in this fixed-stack
demo surface.

The source wave is then generated from the decoded source coefficients by evaluating the
source transfer function
\begin{align}
  H_{\mathrm{src}}(\mathbf{k};\mathbf{u}),
\end{align}
normalizing it in reciprocal space,
\begin{align}
  \widehat{\psi}_{\mathrm{src}}(\mathbf{k};\mathbf{u})
  =
  \frac{H_{\mathrm{src}}(\mathbf{k};\mathbf{u})}
       {\left\lVert H_{\mathrm{src}}(\cdot;\mathbf{u}) \right\rVert_2},
\end{align}
and applying the inverse Fourier transform,
\begin{align}
  \psi_{\mathrm{src}}(\mathbf{r};\mathbf{u})
  = \mathcal{F}^{-1}\!\left\{
      \widehat{\psi}_{\mathrm{src}}(\mathbf{k};\mathbf{u})
    \right\}.
\end{align}
The current implementation then shifts this illumination to the center of the real-space
grid before propagation through the specimen.
For the validated fixed-stack example, the source semiangle cutoff is now chosen large
enough that the source aperture spans multiple reciprocal pixels on the simulation
grid.  In the current example this means a `5.0 mrad` source aperture rather than the
earlier effectively single-pixel regime, so condenser-side controls now produce a
nontrivial pre-specimen wave and a genuinely source-sensitive exit-wave objective.

For a fixed specimen potential stack $\{V_n\}_{n=1}^{N_s}$, the exit wave is
\begin{align}
  \psi_{\mathrm{exit}}(\mathbf{r};\mathbf{u})
  =
  \mathcal{M}_{\Theta_{\mathrm{ms}},\Theta_{\mathrm{pot}}}
  \bigl[\psi_{\mathrm{src}}(\mathbf{r};\mathbf{u})\bigr].
\end{align}
The image-plane wave and image intensity are then formed with the decoded image-side CTF,
\begin{align}
  \psi_{\mathrm{img}}(\mathbf{r};\mathbf{u})
  &=
  \mathcal{F}^{-1}\!\left\{
    H_{\mathrm{img}}(\mathbf{k};\mathbf{u})\,
    \mathcal{F}\{\psi_{\mathrm{exit}}(\mathbf{r};\mathbf{u})\}
  \right\}, \\
  I_{\mathrm{img}}(\mathbf{r};\mathbf{u})
  &= \left|\psi_{\mathrm{img}}(\mathbf{r};\mathbf{u})\right|^2,
\end{align}
and the momentum-plane detector output is
\begin{align}
  I_{\mathrm{diff}}(\mathbf{k};\mathbf{u})
  = \left|\mathcal{F}\{\psi_{\mathrm{exit}}(\mathbf{r};\mathbf{u})\}\right|^2.
\end{align}
The branch therefore emits the tuple
\begin{align}
  y(\mathbf{u}) =
  \left(
    \psi_{\mathrm{src}}(\mathbf{u}),
    I_{\mathrm{img}}(\mathbf{u}),
    I_{\mathrm{diff}}(\mathbf{u}),
    \psi_{\mathrm{exit}}(\mathbf{u})
  \right),
\end{align}
so the pre-specimen source wave, reduced detector data, and the full complex exit wave
remain available to downstream optimization pipelines.
For a batched input tensor $U$, the implementation applies the same forward model to
each control slice independently and preserves the leading batch shape:
\begin{align}
  y(U)
  \in
  \mathbb{C}^{B_1\times\cdots\times B_m\times N_y\times N_x}
\end{align}
for the complex exit-wave branch, with the corresponding real-valued detector outputs
carrying the same leading batch dimensions.

For the current fitting example, the primary optimization target is image contrast
rather than raw intensity:
\begin{align}
  I_{\mathrm{contrast}}(\mathbf{r};\mathbf{u})
  =
  \frac{I_{\mathrm{img}}(\mathbf{r};\mathbf{u})}
       {\overline{I_{\mathrm{img}}^{\star}(\cdot)}}
  - 1,
\end{align}
and the same branch also supports a normalized diffraction metric
\begin{align}
  D_{\mathrm{norm}}(\mathbf{k};\mathbf{u})
  =
  \frac{D(\mathbf{k};\mathbf{u})}
       {\sum_{\mathbf{k}'} D(\mathbf{k}';\mathbf{u})}.
\end{align}
With diffraction weight \(\lambda_D \ge 0\), the demo loss becomes
\begin{align}
  \mathcal{L}(\mathbf{u})
  =
  \frac{1}{N_I}
  \sum_{\mathbf{r}}
  \left(
    I_{\mathrm{contrast}}(\mathbf{r};\mathbf{u})
    -
    I_{\mathrm{contrast}}^{\star}(\mathbf{r})
  \right)^2
  +
  \lambda_D
  \frac{1}{N_D}
  \sum_{\mathbf{k}}
  \left(
    D_{\mathrm{norm}}(\mathbf{k};\mathbf{u})
    -
    D_{\mathrm{norm}}^{\star}(\mathbf{k})
  \right)^2.
\end{align}
This should be interpreted as a differentiable inverse-tuning demonstration, not as a
claim of unique parameter recovery for the current overparameterized control map.
For controls that mainly act in the image-forming branch after the specimen exit wave,
the diffraction contribution can remain nearly constant.  Reciprocal-space fitting is
more informative when the optimized control subset includes source-side parameters that
change the exit-wave illumination itself.  The current validated HRTEM fitting example
therefore screens a mixed source/image plane \((I_{c1}, f_o)\) rather than a purely
post-specimen image-side control plane, while the currently informative live-validated
exit-wave scoring plane is the source-side pair \((I_{c1}, I_{c2})\).
Using the target-image mean as the shared contrast reference also makes the score
surface less sensitive to tiny mean-level numerical differences between otherwise
matching low-contrast images.

For coarse candidate screening before gradient refinement, the same branch now also
supports batched candidate scoring.  Given a candidate tensor
\begin{align}
  U \in \mathbb{R}^{B\times 12}
\end{align}
and a target image $I^\star(\mathbf{r})$, the simulator evaluates all candidates and
forms the batched image metric
\begin{align}
  M_b(\mathbf{r}) &=
  \begin{cases}
    I_{\mathrm{contrast}}(\mathbf{r}; U_b), & \text{if contrast scoring is enabled}, \\
    I_{\mathrm{img}}(\mathbf{r}; U_b), & \text{otherwise},
  \end{cases}
\end{align}
with corresponding target metric
\begin{align}
  M^\star(\mathbf{r}) &=
  \begin{cases}
    I_{\mathrm{contrast}}^\star(\mathbf{r}), & \text{if contrast scoring is enabled}, \\
    I^\star(\mathbf{r}), & \text{otherwise}.
  \end{cases}
\end{align}
with the target contrast defined using the same reference mean
\(
  \overline{I_{\mathrm{img}}^{\star}(\cdot)}
\)
for both the target and all candidate predictions.
If the inverse target also includes the full complex exit wave
\(
  \psi^\star_{\mathrm{exit}}(\mathbf{r})
\),
the simulator can first apply a global phase alignment
\begin{align}
  \widehat{\psi}_{b,\mathrm{exit}}(\mathbf{r})
  &=
  \exp\!\left(
    -i \arg \sum_{\mathbf{r}'}
    \overline{\psi^\star_{\mathrm{exit}}(\mathbf{r}')}
    \psi_{b,\mathrm{exit}}(\mathbf{r}')
  \right)
  \psi_{b,\mathrm{exit}}(\mathbf{r}),
\end{align}
and then form the exit-wave residual
\begin{align}
  L_{\psi,b}
  =
  \frac{1}{N_\psi}
  \sum_{\mathbf{r}}
  \left|
    \widehat{\psi}_{b,\mathrm{exit}}(\mathbf{r})
    -
    \psi^\star_{\mathrm{exit}}(\mathbf{r})
  \right|^2.
\end{align}
If the inverse target instead or additionally includes the pre-specimen source wave
\(
  \psi^\star_{\mathrm{src}}(\mathbf{r})
\),
the simulator can apply the same global phase-alignment convention
\begin{align}
  \widehat{\psi}_{b,\mathrm{src}}(\mathbf{r})
  &=
  \exp\!\left(
    -i \arg \sum_{\mathbf{r}'}
    \overline{\psi^\star_{\mathrm{src}}(\mathbf{r}')}
    \psi_{b,\mathrm{src}}(\mathbf{r}')
  \right)
  \psi_{b,\mathrm{src}}(\mathbf{r}),
\end{align}
and then form the source-wave residual
\begin{align}
  L_{\mathrm{src},b}
  =
  \frac{1}{N_\psi}
  \sum_{\mathbf{r}}
  \left|
    \widehat{\psi}_{b,\mathrm{src}}(\mathbf{r})
    -
    \psi^\star_{\mathrm{src}}(\mathbf{r})
  \right|^2.
\end{align}
The candidate score is then the per-candidate mean squared residual, optionally with
the same diffraction term, the exit-wave term, and the source-wave term,
\begin{align}
  s_b
  =
  \frac{1}{N_I}\sum_{\mathbf{r}} \left(M_b(\mathbf{r}) - M^\star(\mathbf{r})\right)^2
  +
  \lambda_D
  \frac{1}{N_D}
  \sum_{\mathbf{k}}
  \left(
    \widetilde{D}_b(\mathbf{k}) - \widetilde{D}^{\star}(\mathbf{k})
  \right)^2.
  +
  \lambda_\psi L_{\psi,b}
  +
  \lambda_{\mathrm{src}} L_{\mathrm{src},b}.
\end{align}
Here \(\widetilde{D}\) denotes either the raw diffraction intensity or the normalized
diffraction metric depending on the selected scoring mode.
This provides a differentiable screening stage that can rank many microscope-state
proposals against one fixed target image before or alongside subsequent local
optimization.
For local refinement after coarse screening, the simulator also exposes an iterative
optimizer surface that repeatedly evaluates the fixed-stack forward model, computes the
same image, diffraction, and optional wavefield losses, and updates the active control
parameters with a Torch optimizer.  The returned result object records the
loss history together with the initial outputs, final outputs, final control vector,
and the per-step exit-wave and source-wave loss histories when those terms are enabled,
so the inverse-tuning path is available as a reusable differentiable API rather than
only as example-script logic.
Above that, the simulator also exposes a coarse-to-fine orchestration surface for
Cartesian named-parameter grids: it screens the candidate grid, selects the best
screened control vector, instantiates a fresh fixed-stack simulator at that screened
state, and then runs the same differentiable local refinement on top of it.  This
promotes the whole grid-screening plus bounded-refinement workflow to a reusable layer
instead of leaving the orchestration inside one example script.
Because the implementation now also accepts named stage blocks directly, the same
candidate screening may be viewed as acting on a stage-parameter family
\begin{align}
  U_{\mathrm{stage}} =
  \left(
    U_{\mathrm{cond}},
    U_{\mathrm{obj}},
    U_{\mathrm{corr}}
  \right),
\end{align}
which is first assembled into the full candidate tensor $U$ and then scored by the
same formula for $s_b$.  In other words, stage-wise candidate screening is not a
different physical model; it is a different parameterization of the same fixed-stack
HRTEM forward map.
The simulator additionally exposes the corresponding best-candidate selector
\begin{align}
  b^\star = \operatorname*{arg\,min}_{b} s_b,
\end{align}
and returns both the index $b^\star$ and the associated control vector
\begin{align}
  \mathbf{u}^\star = U_{b^\star},
\end{align}
together with the image, diffraction, and exit-wave outputs for that best-scoring
candidate.  This yields a direct coarse-to-fine workflow: screen a batched family of
microscope-state proposals, extract $\mathbf{u}^\star$, and continue refinement from
that candidate by gradient-based optimization if needed.
For workflows that keep several promising candidates, the same score tensor may be
ranked by
\begin{align}
  b_1^\star, \dots, b_k^\star
  = \operatorname*{arg\,topk\,min}_{b} s_b,
\end{align}
with associated vectors
\begin{align}
  \mathbf{u}_{(1)}^\star, \dots, \mathbf{u}_{(k)}^\star,
\end{align}
returned in ascending score order.  This allows downstream optimization to branch from
the best few microscope-state proposals instead of only the single global minimum over
the screened batch.
When the candidates come from a Cartesian parameter scan such as
\begin{align}
  I_{c1} \in \{I_{c1}^{(1)}, \dots, I_{c1}^{(M)}\},
  \qquad
  f_o \in \{f_o^{(1)}, \dots, f_o^{(N)}\},
\end{align}
the implementation first builds the corresponding parameter grid tensor and then applies
the same score map $s_b$, best-candidate selector $b^\star$, and top-$k$ ranking
operators.  In other words, grid screening is the same coarse-search primitive as the
arbitrary batched candidate case, but with a structured Cartesian parameterization of
the proposal family.

\end{document}
